\documentclass[11pt,a4paper,sans]{moderncv} % Font sizes: 10, 11, or 12; paper sizes: a4paper, letterpaper, a5paper, legalpaper, executivepaper or landscape; font families: sans or roman



\moderncvstyle{classic} % CV theme - options include: 'casual' (default), 'classic', 'oldstyle' and 'banking'
\moderncvcolor{blue} % CV color - options include: 'blue' (default), 'orange', 'green', 'red', 'purple', 'grey' and 'black'

\usepackage{lipsum} % Used for inserting dummy 'Lorem ipsum' text into the template

\usepackage[scale=0.8]{geometry} % Reduce document margins
%\setlength{\hintscolumnwidth}{3cm} % Uncomment to change the width of the dates column
%\setlength{\makecvtitlenamewidth}{10cm} % For the 'classic' style, uncomment to adjust the width of the space allocated to your name

%----------------------------------------------------------------------------------------
%	NAME AND CONTACT INFORMATION SECTION
%----------------------------------------------------------------------------------------

\firstname{Soumyabrata} % Your first name
\familyname{Pal} % Your last name

% All information in this block is optional, comment out any lines you don't need
\title{Teaching Statement}
%\address{W. Ethan Eagle}{}
\mobile{+91-7604039429}
%\phone{(000) 111 1112}
%\fax{(000) 111 1113}
\email{soumyabratapal13@gmail.com}
\homepage{https://soumyabratap.github.io/}% The first argument is the url for the clickable link, the second argument is the url displayed in the template - this allows special characters to be displayed such as the tilde in this example
%\extrainfo{additional information}
%\photo[70pt][0.4pt]{pictures/picture} % The first bracket is the picture height, the second is the thickness of the frame around the picture (0pt for no frame)
%\quote{"A witty and playful quotation" - John Smith}

%----------------------------------------------------------------------------------------

\begin{document}
\makecvtitle % Print the CV title
I am fortunate to have experienced a diverse set of teaching styles through the course of my graduate and undergraduate studies. My teachers and mentors played a crucial role in shaping my research and career goals. I aspire to have a meaningful impact on the next generation through my enthusiasm and passion for teaching. My teaching style, outlined in the next section, has been formed through the course of several experiences as a student and researcher. These include assisting a graduate mathematics and computer science course, giving well-received and well-appreciated guest lectures at different courses, and through countless presentations at seminars and reading groups.

\section{Teaching Principles}

$\bullet{}$ \textbf{Developing intuition through visualizations:} I believe that the primary goal of teaching is to foster intuition and curiosity. Once students develop the right intuition, explaining the most profound concepts becomes a relatively easy task. As a teaching assistant for Combinatorics and Graph Theory, I was faced with the challenge of teaching some core constructs to undergraduate and graduate students with no prior knowledge of proof techniques. I realized that visualizations are a convenient tool to nurture insights. It allowed me to present complex ideas as a sequence of multiple intuitive modules with natural transitions between them. I noted that this strategy not only enhanced the students' understanding, but also improved their retention.

$\bullet{}$\textbf{Team-based learning:} In my opinion, creativity is best promoted in question-friendly environments and small groups. Setting up a comfortable and interactive atmosphere in classrooms would allow students to speak without the fear of failure. For large classes, having group assignments and interacting with the students in smaller help-sessions would achieve this. Many students appreciated this approach that I took as a teaching assistant.

$\bullet{}$ \textbf{Example based teaching:} Another strategy I have found useful is to use exciting puzzles to motivate complicated algorithmic paradigms. This methodology was employed by Professor David Barrington to teach the Combinatorics and Graph Theory course at Umass Amherst. It helped the students generate new ideas and think about the problems intuitively. Moreover, the diverse set of ideas generated in the classroom also aided lateral thinking. 

$\bullet{}$ \textbf{Having structured courses:} My duties as a TA  included assisting with curriculum development and development of the course material. Through these experiences, I realized the importance of well-structured courses and stimulating assignments. The students appreciated having a systematic course that builds up the theories from past lectures. I would also encourage regular evaluations through exams, quizzes, and assignments. These assignments should be challenging enough to inculcate a creative thought process while applying the concepts taught in class.

$\bullet{}$ \textbf{Encouraging diversity and community outreach:} In a classroom, as in research, diversity in ethnic, gender, and academic backgrounds is crucial to encourage diverse perspectives and learn from diverse experiences. As a young faculty, I will direct my efforts to retain diversity through an inclusive classroom environment, which is as essential as diverse recruiting. As a teacher involved in research, I will engage in outreach activities aimed at explaining the benefits and advances of research in my field to a broader audience. Since every student is different, a particular teaching style may not benefit everyone equally. Through regular feedback, I aim to be flexible and adaptive in my teaching approaches.

\newpage 

\section{Teaching Expertise}

My research interests span a broad range of topics in Mathematics, Computer Science and Machine Learning. I believe I am qualified to teach courses in related areas, at all levels of the undergraduate curriculum as well as advanced graduate levels. I can teach undergraduate courses on Algorithms, Discrete Mathematics, and Theory of Computation. At the graduate level, I can teach advanced courses on Algorithms, Foundations of Data Science and Machine Learning, and Optimization (Convex and Non-convex). Other special topics I would be extremely interested in teaching are courses on:


$\bullet{}$ \textbf{Machine Learning Theory:} This course will cover many recent techniques applied for statistical reconstruction problems such as tensor based algorithms, spectral algorithms, non-negative matrix factorization, non-convex optimization, etc.


$\bullet{}$\textbf{Transformers and Generative AI} This course would cover the recent methods and techniques in Large Model that are based on Transformers. This course will be particularly relevant to industry and will cover topics on a variety of approaches for latency, memory and computational trade-offs.


$\bullet{}$\textbf{Mathematical Toolkit:}  This course would introduce the mathematical tools that are commonly used in other advanced areas. I feel many students interested in the foundations of data science and machine learning would benefit from such a course. \newline

Moreover, seminars and reading groups are an effective and efficient way to introduce a variety of topics to students. I want to continue organizing such seminars and reading groups regularly in the future. I look forward to continually learn from my past and future experiences and grow as a teacher and a mentor.


\end{document}

